% Source: https://artofproblemsolving.com/community/c3h1148483_expected_value
% Stable identity: aops-topic-1148483-q01
A coin shows tails only if Kelvin flipped it to tails and then hopped left off it. Follow the process:
\begin{itemize}
\item While he keeps flipping heads, he hops right onto pads that have never been touched (heads), so he cannot fall asleep.
\item At his first tails he hops left onto a pad showing heads (either the pad he just came from, which he left showing heads, or an untouched pad), and the pad to his right now shows tails.
\item From then on, every tails moves him one more pad to the left onto a heads pad (again leaving a tails behind him), while his first heads moves him right onto the pad he just left, which shows tails: he falls asleep.
\end{itemize}
So the time is (number of flips up to and including the first tails) plus (number of further flips up to and including the first heads). Each is a geometric random variable with success probability $\dfrac{1}{2}$, with mean $2$. The expected time is $2+2=\boxed{4}$ minutes.
